\documentclass[11pt]{article}

% ---------------------------------------------------------------------------
%  The Cauldron — Full-Protocol Security Audit (from-scratch review)
%  Compile:  pdflatex CauldronSecurityAudit.tex   (run twice for refs/TOC)
% ---------------------------------------------------------------------------

\usepackage[a4paper,margin=1in]{geometry}
\usepackage{xcolor}
\usepackage{listings}
\usepackage{booktabs}
\usepackage{amsmath,amssymb}
\usepackage{enumitem}
\usepackage{titlesec}
\usepackage{fancyhdr}
\usepackage{tcolorbox}
\usepackage{longtable}
\usepackage[hidelinks]{hyperref}
\usepackage{parskip}

% --- palette --------------------------------------------------------------
\definecolor{void}{HTML}{0E0A1A}
\definecolor{lime}{HTML}{6E8B00}
\definecolor{violet}{HTML}{5E43C9}
\definecolor{sevCrit}{HTML}{B00020}
\definecolor{sevHigh}{HTML}{E4572E}
\definecolor{sevMed}{HTML}{C08B00}
\definecolor{sevLow}{HTML}{3A7D44}
\definecolor{sevInfo}{HTML}{5A6472}
\definecolor{codebg}{HTML}{F5F4FA}

\hypersetup{colorlinks=true,linkcolor=violet,urlcolor=violet,citecolor=violet}

% --- listings -------------------------------------------------------------
\lstdefinelanguage{Solidity}{
  keywords={pragma,contract,interface,library,function,modifier,event,error,
    struct,mapping,address,uint256,uint128,uint160,uint96,uint48,uint16,uint8,
    int24,int128,bool,bytes,bytes32,string,memory,storage,calldata,external,
    public,internal,private,view,pure,payable,returns,return,if,else,for,while,
    require,revert,emit,new,immutable,constant,override,virtual,using,is,
    unchecked,type,msg,block,this,true,false},
  keywordstyle=\color{violet}\bfseries,
  comment=[l]{//}, morecomment=[s]{/*}{*/},
  commentstyle=\color{sevInfo}\itshape,
  stringstyle=\color{sevLow}, morestring=[b]",
  sensitive=true,
}
\lstset{
  language=Solidity, backgroundcolor=\color{codebg},
  basicstyle=\ttfamily\footnotesize, breaklines=true,
  frame=single, rulecolor=\color{codebg!60!black},
  numbers=left, numberstyle=\tiny\color{sevInfo}, numbersep=8pt,
  showstringspaces=false, xleftmargin=2.2em, framexleftmargin=2em, aboveskip=1em,
}

% --- headers --------------------------------------------------------------
\pagestyle{fancy}\fancyhf{}
\lhead{\small\textsc{The Cauldron — Full-Protocol Audit}}
\rhead{\small\thepage}
\renewcommand{\headrulewidth}{0.4pt}

\titleformat{\section}{\Large\bfseries\color{void}}{\thesection}{0.6em}{}
\titleformat{\subsection}{\large\bfseries\color{violet}}{\thesubsection}{0.6em}{}

% --- finding box ----------------------------------------------------------
\newtcolorbox{finding}[2]{colback=codebg,colframe=#1,boxrule=1pt,
  arc=2pt,left=8pt,right=8pt,top=6pt,bottom=6pt,
  title={\bfseries #2},fonttitle=\color{white},coltitle=white}

\newcommand{\sev}[2]{\textcolor{#1}{\textbf{#2}}}
\newcommand{\addr}[1]{\texttt{\footnotesize #1}}

% ==========================================================================
\begin{document}

\begin{titlepage}
\centering
\vspace*{2.2cm}
{\Huge\bfseries\color{void} The Cauldron}\\[0.35cm]
{\LARGE\color{violet} Full-Protocol Security Audit}\\[0.2cm]
{\large\itshape a from-scratch review of the eternal-token machine}\\[1.4cm]

\begin{tcolorbox}[colback=codebg,colframe=violet,width=0.8\textwidth,arc=3pt]
\centering
\textbf{Subject:} MagicFrens ``Cauldron'' — Uniswap~v4 hook-native, self-relaunching
token machine, with the new \emph{prime-buy}, \emph{proposer-incentive} and
\emph{V2-upgrade (modular controller)} subsystems.\\[0.4cm]
\textbf{Review type:} Full re-audit (all findings re-derived; no prior report reused).\\[0.2cm]
\textbf{Network:} Ethereum Sepolia (test) $\rightarrow$ mainnet-candidate.\\[0.2cm]
\textbf{Date:} 31~August~2026.
\end{tcolorbox}

\vfill
{\small\color{sevInfo} Prepared as an internal pre-deployment review. Not a warranty.
A smart contract carries residual risk regardless of review depth.}
\end{titlepage}

\tableofcontents
\newpage

% ==========================================================================
\section{Abstract}

The Cauldron is a single, permissionless, self-perpetuating token machine built as a
Uniswap~v4 hook. One ``iteration'' lives as a v4 pool whose entire supply is held as
liquidity: a thin \emph{active} band that sets price, and a large \emph{out-of-range
reserve} that backs a redemption floor and funds 1:1 migration. When an iteration's
rolling volume dies, anyone may \texttt{relaunch()}: the dead LP is recovered, its tokens
burned, and a fresh iteration is born from a governor-selected community proposal — seeded
with a real first-block ``green-candle'' market buy. Genesis MiFrens holders earn a
permanent founder's cut of every iteration; per-collection ``legacy'' floors and an
in-hook perpetuals engine layer on top.

This review is a \emph{from-scratch} re-derivation of the protocol's security posture,
undertaken before the next deployment (``r30''). It covers the full contract set and, in
particular, three subsystems added since the last review:

\begin{enumerate}[leftmargin=1.4em]
  \item \textbf{Prime buy} — an owner-funded, personal-ETH first-block market buy at
        genesis whose output goes to the treasury for a later OG airdrop (net demand,
        zero supply dilution).
  \item \textbf{Proposer incentive} — a tiny per-iteration fee slice that rewards whoever
        authored the winning proposal (a decentralization flywheel).
  \item \textbf{Modular V2-upgrade path} — a governed controller handoff that lets a
        future ``cellular division'' V2 inherit the live liquidity \emph{by ownership
        transfer, without a teardown}, protected by a hard exit guarantee and a guardian
        veto.
\end{enumerate}

One \sev{sevMed}{Medium} issue was found \emph{and fixed during this review}
(A-01: an untrusted external call in the swap fee hot path). No unresolved Critical or
High severity issue that risks principal was identified. The dominant residual risk is
\emph{governance centralization}, which the protocol addresses by an explicit,
deliberately-chosen ``powerful but telegraphed, with a guaranteed floor exit'' trust
model rather than by hard on-chain constraints (\S\ref{sec:trust}).

% ==========================================================================
\section{Methodology}

The review followed a manual, adversarial reading of every state-changing external entry
point, cross-referenced against automated size/compile checks and the project's own
Foundry fork-test suite (141 tests against live Sepolia v4). For each entry point we asked:
who may call it; what value can move; is Checks-Effects-Interactions respected; is any
external call made to an untrusted address; and can a swap, relaunch, or redemption be
bricked or front-run. Findings were classified by the standard matrix:

\begin{center}\small
\begin{tabular}{ll}
\toprule
\sev{sevCrit}{Critical} & Direct principal loss, no special conditions.\\
\sev{sevHigh}{High}     & Principal loss under specific but reachable conditions.\\
\sev{sevMed}{Medium}    & Limited loss, griefing, or a governance/centralization lever.\\
\sev{sevLow}{Low}       & Minor / best-practice / defence-in-depth.\\
\sev{sevInfo}{Info}     & Observation, or a positive design confirmation.\\
\bottomrule
\end{tabular}
\end{center}

% ==========================================================================
\section{Scope}

\begin{center}\small
\begin{longtable}{p{0.30\textwidth}p{0.62\textwidth}}
\toprule
\textbf{Contract} & \textbf{Role}\\
\midrule
\endhead
\texttt{CauldronRegistry} & Orchestrator: summon, relaunch, redemption floor, prime buy,
  emergency \& V2-upgrade governance. Custody of the live position-NFTs.\\
\texttt{CauldronHook} & v4 hook: swap fees, anti-snipe surtax, volume/death, legacy
  buyback, guild \& proposer routing, gacha, perp-fee routing.\\
\texttt{PoolOps} (library) & Delegatecalled LP primitives: seed, green-candle buy,
  prime buy, reserve add/claim, migration, \texttt{CauldronToken} CREATE2 deploy.\\
\texttt{CauldronToken} & Fixed-supply, non-mintable, non-freezable ERC-20 per iteration.\\
\texttt{PerpEngine} / \texttt{PerpVault} & Hook-native perpetuals + community PLV vault.\\
\texttt{CollectionLedger} & Per-collection legacy-floor cap table (accounting only).\\
\texttt{MiFrensGenesis} / \texttt{MiFrensDividend} & OG NFT + founder fee dividend.\\
\texttt{CauldronGovernor} / \texttt{Factory} / \texttt{GachaRouter} & Proposals,
  collection deploy, crystal gacha router.\\
\texttt{ReserveLib} / \texttt{LaunchLib} & Pure tick / naming helpers.\\
\bottomrule
\end{longtable}
\end{center}

\noindent Timelock: OpenZeppelin \texttt{TimelockController} is the registry's
\texttt{emergencyAdmin} and the owner of the hook \& perp engine post-deploy.

% ==========================================================================
\section{Access-Control \& Fund-Mover Matrix}

Every function that can move ETH or tokens \emph{out} of the system, and its gate:

\begin{center}\small
\begin{longtable}{p{0.33\textwidth}p{0.29\textwidth}p{0.30\textwidth}}
\toprule
\textbf{Fund-mover} & \textbf{Caller gate} & \textbf{Extra guard}\\
\midrule
\endhead
\texttt{emergencyWithdrawLP} & emergencyAdmin (timelock) & armed+delay, exit-open, vetoable, nonReentrant\\
\texttt{emergencySweep} & emergencyAdmin (timelock) & armed+delay, vetoable, nonReentrant\\
\texttt{migrateToSuccessor} & emergencyAdmin (timelock) & armed+delay, exit-open, vetoable, nonReentrant\\
\texttt{redeemFren / claimByBurn} & any NFT owner / holder & bounded by floor; paid from reserve LP\\
\texttt{buyTreasuryFren / recycle} & permissionless & pays/pulls at the accounted floor\\
\texttt{sweepPrimeBuy} & primeFunder & CEI (balance zeroed first)\\
\texttt{releaseRelaunchETH} & registry only & sends tracked \texttt{relaunchETH} var\\
\texttt{claimProposerFees} & self (pull) & CEI + nonReentrant\\
\texttt{PerpVault / PerpEngine.withdrawPlvTo} & vault / share-owner & onlyVault, notNested, nonReentrant\\
\bottomrule
\end{longtable}
\end{center}

\noindent No path transfers principal to an arbitrary, caller-supplied address without
either (a) a timelock+veto governance gate or (b) a floor-bounded, permissionless
redemption. The one place an attacker-controlled address received a push — the proposer
slice — was converted to a pull during this review (A-01).

% ==========================================================================
\section{Findings Summary}

\begin{center}\small
\begin{longtable}{llp{0.42\textwidth}l}
\toprule
\textbf{ID} & \textbf{Severity} & \textbf{Title} & \textbf{Status}\\
\midrule
\endhead
A-01 & \sev{sevMed}{Medium}  & Untrusted external call in the swap fee hot path (proposer push) & \textbf{Fixed}\\
A-02 & \sev{sevMed}{Medium}  & Governance centralization (custody \& upgrade levers) & Mitigated (by design)\\
A-03 & \sev{sevLow}{Low}     & Hook re-point exposes the \texttt{relaunchETH} buffer & Acknowledged\\
A-04 & \sev{sevLow}{Low}     & Unchecked ERC-20 return values on protocol tokens & Acknowledged\\
A-05 & \sev{sevInfo}{Info}   & Stale position pointers after \texttt{migrateToSuccessor} & By design\\
A-06 & \sev{sevInfo}{Info}   & Prime-buy recipient-zero fail-safe & By design\\
A-07 & \sev{sevInfo}{Positive} & Successor handoff preserves price \& liquidity (no teardown) & Confirmed\\
A-08 & \sev{sevInfo}{Positive} & Exit guarantee: redemptions forced open while armed & Confirmed\\
A-09 & \sev{sevInfo}{Positive} & Reserve solvency invariant (fuzz + fork proven) & Confirmed\\
\bottomrule
\end{longtable}
\end{center}

% ==========================================================================
\section{Detailed Findings}

\subsection{A-01 --- Untrusted external call in the swap fee hot path \sev{sevMed}{(Medium, Fixed)}}
\begin{finding}{sevMed}{A-01 · CauldronHook · Reentrancy / Untrusted Call}
The proposer-incentive slice was originally \emph{pushed} to
\texttt{activeProposer} with a value-bearing \texttt{.call} inside
\texttt{\_routeEthFee}, which runs during \texttt{beforeSwap}/\texttt{afterSwap}.
\end{finding}

\noindent\textbf{Impact.} \texttt{activeProposer} is attacker-controllable: anyone may
submit a proposal, so a winning proposer can point the slice at a contract of their
choosing. A value push to that address inside the un-guarded hook swap entrypoint is a
classic reentrancy surface — the recipient executes mid-swap, before the fee's
guild/floor/relaunch state updates complete, and could attempt to re-enter the hook,
registry, or a nested swap.

\noindent\textbf{Resolution (this review).} Converted to a \textbf{pull pattern}. The
slice now accrues to \texttt{proposerOwed[proposer]} (ETH retained in the hook, tracked by
variable exactly like \texttt{relaunchETH}/\texttt{legacyBuffer}), and the proposer later
withdraws via \texttt{claimProposerFees()} (CEI + \texttt{nonReentrant}). The swap hot
path now performs \emph{zero} external calls to untrusted addresses.

\begin{lstlisting}
// before (push, untrusted call mid-swap):
(bool okP,) = prop.call{value: wantProp}("");
// after (pull, no call in the hot path):
proposerOwed[prop] += wantProp; feeAmount -= wantProp; emit ProposerFunded(prop, wantProp);
\end{lstlisting}

\noindent Verified by \texttt{test\_ProposerIncentive\_PaysProposer\_OnFork} (accrues,
then claims exactly, then zeroes).

\subsection{A-02 --- Governance centralization of custody \& upgrade levers \sev{sevMed}{(Medium, mitigated by design)}}
\label{sec:trust}
\begin{finding}{sevMed}{A-02 · CauldronRegistry / CauldronHook · Centralization}
\texttt{emergencyWithdrawLP}, \texttt{emergencySweep}, \texttt{migrateToSuccessor} and the
hook's \texttt{setRegistryOverride} let the governance timelock move the live liquidity or
re-point the controller. This is the protocol's largest trust surface.
\end{finding}

\noindent\textbf{Design posture.} The project deliberately chooses a \emph{``powerful but
telegraphed, with a guaranteed exit''} model over hard ``cannot-rug-by-construction''
constraints, in order to keep a break-glass escape and to leave a future V2 unconstrained.
The residual risk is bounded by \textbf{five independent layers}:

\begin{enumerate}[leftmargin=1.4em]
  \item \textbf{Timelock delay} — every custody action must be \texttt{armEmergency}'d and
        wait \texttt{emergencyDelay} (mainnet: days), with a public event.
  \item \textbf{Forced-open exit} — the instant any action is armed,
        \texttt{\_redeemBlocked()} returns false regardless of \texttt{redemptionPaused};
        holders can always redeem at floor \emph{before} anything moves. This closes the
        ``pause-then-drain'' trap.
  \item \textbf{Guardian veto} — a block-only \texttt{guardian} may
        \texttt{vetoEmergency()} to cancel an armed action; it can never propose or move
        funds.
  \item \textbf{Reserve solvency} — the exit is \emph{funded}: the out-of-range reserve is
        intact during the whole window (the withdraw is the last step), so a bank-run at
        floor is payable (A-09).
  \item \textbf{Community ratification (recommended for mainnet)} — gate
        \texttt{setSuccessor} behind a genesis-fren vote so one key cannot hand off.
\end{enumerate}

\noindent\textbf{Residual.} Layers 1--4 protect \emph{attentive} holders and passive
holders alike (the veto is passive-safe). A fully-compromised timelock \emph{and} guardian
could still, after the delay, execute a withdraw — but holders retained a funded floor
exit throughout. Recommendation: a multi-sig guardian distinct from the timelock proposer
set, and a $\geq$7-day \texttt{emergencyDelay} on mainnet.

\subsection{A-03 --- Hook re-point exposes the \texttt{relaunchETH} buffer \sev{sevLow}{(Low)}}
\begin{finding}{sevLow}{A-03 · CauldronHook · Access Control}
\texttt{setRegistryOverride} (onlyOwner = timelock) re-points the hook's trusted
\texttt{registry}. A malicious re-point could let a substitute ``registry'' call
\texttt{releaseRelaunchETH} and drain the accrued relaunch buffer.
\end{finding}

\noindent\textbf{Impact.} Bounded to the \texttt{relaunchETH} buffer (fees accrued since
the last relaunch) — \emph{not} the LP, which is owned by the registry address, not the
hook. Requires a compromised/malicious timelock (same base trust as A-02).
\textbf{Recommendation.} Perform the hook re-point in the \emph{same} governance ceremony
as \texttt{migrateToSuccessor} (which carries the exit window), and treat it as a
custody-class action operationally.

\subsection{A-04 --- Unchecked ERC-20 return values \sev{sevLow}{(Low)}}
\begin{finding}{sevLow}{A-04 · Multiple · Unchecked Calls}
Several \texttt{IERC20.transfer/transferFrom} calls do not check the boolean return
(e.g. registry airdrop/sweep, perp token funding).
\end{finding}

\noindent\textbf{Impact.} Low: the only tokens handled are the protocol's own
\texttt{CauldronToken} (OpenZeppelin-based, reverts on failure) and ETH. A non-standard
token is never in scope. \textbf{Recommendation.} Adopt \texttt{SafeERC20} uniformly before
mainnet as defence-in-depth against any future integration of a non-reverting token.

\subsection{A-05 --- Stale position pointers after migration \sev{sevInfo}{(Info, by design)}}
\begin{finding}{sevInfo}{A-05 · CauldronRegistry · Logic}
After \texttt{migrateToSuccessor}, \texttt{generationPositionId}/\texttt{...ReservePositionId}
still reference NFTs the registry no longer owns; a subsequent \texttt{redeemFren} on the
old registry would revert.
\end{finding}

\noindent\textbf{Rationale.} Intended: the old registry is \emph{defunct} post-handoff;
all custody and user flows move to the successor. Reverting (rather than mis-paying) is the
fail-safe outcome. Front-ends must route redemptions to the successor after a migration.

\subsection{A-06 --- Prime-buy recipient-zero fail-safe \sev{sevInfo}{(Info, by design)}}
\begin{finding}{sevInfo}{A-06 · CauldronRegistry / PoolOps · Input Validation}
The genesis prime buy sends bought tokens to
\texttt{airdropWallet}, falling back to \texttt{primeFunder}. If both were zero the
\texttt{take} to \texttt{address(0)} would revert.
\end{finding}

\noindent\textbf{Rationale.} \texttt{primeBuyEth} can only be funded by
\texttt{primeFunder}, which is therefore non-zero whenever a prime buy exists — the
recipient is non-zero by construction. A misconfiguration reverts the summon
\emph{loudly} rather than burning tokens. Acceptable fail-safe.

\subsection{A-07 --- Successor handoff preserves price \& liquidity \sev{sevInfo}{(Positive)}}
\begin{finding}{sevLow}{A-07 · CauldronRegistry · Upgradeability}
\texttt{migrateToSuccessor} transfers the position-NFTs' \emph{ownership}; it never
withdraws liquidity.
\end{finding}

\noindent\textbf{Confirmation.} Uniswap~v4 positions are ERC-721s; transferring the token
changes only the owner-of-record, leaving the pool's \texttt{sqrtPrice}, tick, and position
liquidity untouched. Fork-proven byte-for-byte by
\texttt{test\_MigrateToSuccessor\_MovesOwnership\_NoTeardown\_OnFork}. This is what makes a
V2 controller a \emph{drop-in swap} rather than an emergency-withdraw + redeploy.

\subsection{A-08 --- Exit guarantee: redemptions forced open while armed \sev{sevInfo}{(Positive)}}
\begin{finding}{sevLow}{A-08 · CauldronRegistry · Safety}
\texttt{\_redeemBlocked() = redemptionPaused \&\& emergencyReadyAt == 0}.
\end{finding}

\noindent\textbf{Confirmation.} The redemption pause cannot be used to trap holders ahead
of a custody move: arming any emergency action forces the floor exit open for the whole
window. Fork-proven by
\texttt{test\_ExitForcedOpenWhileArmed\_AndGuardianVeto\_OnFork} (paused$\rightarrow$blocked;
armed$\rightarrow$open; vetoed$\rightarrow$blocked again; non-guardian cannot veto).

\subsection{A-09 --- Reserve solvency invariant \sev{sevInfo}{(Positive)}}
\begin{finding}{sevLow}{A-09 · CauldronRegistry / CollectionLedger · Economic}
The out-of-range reserve always covers 1:1 migration + unclaimed genesis + crystallized
legacy entitlements.
\end{finding}

\noindent\textbf{Confirmation.} The relaunch reserve is sized as
\texttt{TOTAL\_SUPPLY - newActive}, where \texttt{newActive} is the survived circulating
supply net of unclaimed genesis and legacy carve. The \texttt{CollectionLedger}
conservation invariant is fuzz-proven over 128{,}000 operations; 1:1 migration out of the
freshly-bought reserve is fork-proven under legacy-carve stress.

% ==========================================================================
\section{Invariants \& Test Coverage}

\begin{itemize}[leftmargin=1.4em]
  \item \textbf{141 / 141} Foundry tests pass (fork against live Sepolia v4).
  \item All in-scope deployable contracts are under the EIP-170 24{,}576-byte limit
        (registry 23.7\,KB, hook 22.8\,KB, engine 24.4\,KB, PoolOps 17.0\,KB).
  \item Key fork proofs: genesis prime buy (treasury receives GNOME, zero dilution, spot
        pumps); proposer accrual/claim; exit-guarantee \& guardian veto; migration
        no-teardown; 1:1 migration solvency; perp auto-migrate at relaunch.
\end{itemize}

% ==========================================================================
\section{Gas Optimization Audit}

A full-scope gas review was run across all contracts, prompted by a concrete failure:
\textbf{in-swap liquidation of a short position silently no-oped under a normal
swap's gas budget}. A short liquidation was measured at \textbf{$\approx$957k gas}; the
hook only forwards \texttt{gasleft $-$ 180k} to the sweep and only fires it above
$\approx$430k, so an ordinary router/UI swap never carried enough gas to settle a
short — the position stayed open despite being underwater at the mark, and no
Liquidatoor badge was struck.

\begin{center}\small
\begin{longtable}{llp{0.40\textwidth}l}
\toprule
\textbf{ID} & \textbf{Contract} & \textbf{Optimization} & \textbf{Impact}\\
\midrule
\endhead
G-01 & PerpEngine & \texttt{OBS\_CARDINALITY} 128$\rightarrow$32 — the TWAP mark
  loop cold-read 128 slots but a 5-min window at 30s spacing can never use $>$$\sim$10 & $\approx$$-$200k on the first mark loop\\
G-02 & PerpEngine & Cache the TWAP mark ONCE per liquidation — it was recomputed
  2--3$\times$ (\texttt{\_underwater} + notional-cap); now the value is computed once
  and reused (\texttt{\_underwaterVal}) & removes 1--2 redundant mark loops\\
G-03 & PerpEngine & Liquidatoor badge is now \textbf{hybrid}: a gas-bounded
  best-effort auto-mint in-swap (matching the creature-NFT UX) that downgrades to a
  claimable credit (\texttt{claimLiquidatorBadges}) only when gas is tight — so a
  liquidation never reverts on the trophy & removes the badge from the critical path\\
G-04 & All & Loop hygiene: cache \texttt{.length}, \texttt{unchecked \{++i\}},
  and cache the \texttt{minted} counter in \texttt{MiFrensGenesis.mint} (was a
  storage read+write every iteration, up to 100$\times$/mint) & per-iteration SLOAD/SSTORE
  savings; also \emph{shrinks} bytecode\\
\bottomrule
\end{longtable}
\end{center}

\noindent\textbf{Result.} A short liquidation dropped from $\approx$957k to
$\approx$450k gas, and the in-swap path (which reuses the already-open PoolManager
lock and no longer mints the badge) is lighter still — so swaps carrying a normal
router/UI gas budget now trigger short liquidations, closing the ``every swap
liquidates'' gap. Ultra-tight swaps still defer to the permissionless keeper
(\texttt{liquidate}/\texttt{sweepLiquidations}), which mainnet MEV bots run. All 141
fork tests remain green and every deployable contract stays under EIP-170 (the
\texttt{unchecked} loops reclaimed bytecode headroom on the byte-tight engine).

\noindent\textbf{Note.} These are code-level improvements; the currently-deployed r30
contracts realise them on the next deployment (r31). No behavioural or economic
change — badge attribution is preserved by the event; the pull-claim mints the same
trophy into the live collection.

% ==========================================================================
\section{Conclusion}

The Cauldron's principal-holding paths are, to the depth of this review, free of
unresolved Critical/High vulnerabilities. The one Medium issue surfaced by this
review — an untrusted external call in the swap fee path — was fixed in place with a pull
pattern. The protocol's remaining material risk is governance centralization, which it
addresses not by removing power but by making every dangerous action slow, visible,
vetoable, and — crucially — always preceded by a \emph{funded, un-blockable floor exit}.
For mainnet we recommend: a $\geq$7-day emergency delay, a guardian multisig independent of
the timelock proposers, \texttt{SafeERC20} throughout, and gating \texttt{setSuccessor}
behind a genesis-fren ratification vote. With those, the ``eternal machine'' can upgrade
itself toward the V2 cellular-division design as a governed pointer swap on live
liquidity — never a teardown.

\end{document}
